\section{Données}\label{sec:donnees}

La base \emph{MAP\_DATA\_WELD} (T.~Cool et H.~K.~D.~H.~Bhadeshia, Cambridge, 1996)
rassemble des mesures tirées de nombreux articles sur le \emph{métal déposé}, le métal
fondu qui forme le cordon de soudure. Elle compte \textbf{1652 lignes et 44 colonnes}.
Une valeur absente est notée \code{N} : non rapportée, et non nulle.

\outrouver{données brutes \code{data/welddb/welddb.data} (à télécharger, voir
\code{README.md}) ; notebook \code{data_preprocessing.ipynb}, « Data Exploration ».}

\subsection{Exploration}\label{sec:exploration}

\subsubsection{Structure}

Les colonnes 1 à 30, choisies avant ou pendant le soudage, sont les entrées. Les
colonnes 31 à 43, résultats d'essais, sont les grandeurs à prédire
(tableau~\ref{tab:colonnes}).

\begin{table}[!htbp]
  \centering\small
  \begin{tabularx}{\linewidth}{@{}l c Y l@{}}
    \toprule
    \textbf{Groupe} & \textbf{Col.} & \textbf{Variables (unité)} & \textbf{Rôle} \\
    \midrule
    Composition       & 1--12  & C, Si, Mn, S, P, Ni, Cr, Mo, V, Cu, Co, W (\% massique) & entrée \\
    Éléments mineurs  & 13--21 & O, Ti, N, Al, B, Nb, Sn, As, Sb (ppm massique) & entrée \\
    Soudage           & 22--28 & courant (\si{\ampere}), tension (\si{\volt}), type de
                                 courant AC/DC, polarité, énergie de soudage
                                 (\si{\kilo\joule\per\milli\metre}), température entre
                                 passes (\si{\celsius}), procédé & entrée \\
    Traitement thermique & 29--30 & température (\si{\celsius}) et durée (\si{\hour}) ;
                                 \code{0/0} = soudure brute & entrée \\
    Propriétés mécaniques & 31--38 & limite d'élasticité, résistance à la traction
                                 (\si{\mega\pascal}), allongement, striction (\%),
                                 température et énergie Charpy (\si{\celsius},
                                 \si{\joule}), dureté, FATT & sortie \\
    Microstructure    & 39--43 & proportions de cinq structures cristallines (\%) & sortie \\
    Identifiant       & 44     & \texttt{Weld ID}, qui indique aussi l'article d'origine & groupe \\
    \bottomrule
  \end{tabularx}
  \caption{Organisation des 44 colonnes de la base.}
  \label{tab:colonnes}
\end{table}

\paragraph{Vocabulaire.}
\begin{itemize}
  \item \emph{ppm} : partie par million en masse (\SI{1}{ppm} = 0,0001\,\%).
  \item \emph{Température entre passes} : température du métal quand on dépose une
    nouvelle couche de soudure.
  \item \emph{Traitement thermique} : réchauffage après soudage pour relâcher les
    contraintes ; il modifie le métal.
  \item \emph{Limite d'élasticité, résistance à la traction} : efforts à partir desquels
    le métal se déforme pour de bon, puis casse.
  \item \emph{Allongement, striction} : capacité à se déformer sans casser (ductilité).
  \item \emph{Essai Charpy} : énergie absorbée par une éprouvette entaillée cassée d'un
    coup de pendule ; plus elle est grande, moins le métal est fragile. Elle dépend
    fortement de la température d'essai.
\end{itemize}

\paragraph{Une ligne correspond à un essai, pas à une soudure.} Plusieurs éprouvettes
d'un même cordon sont testées différemment. Le cordon \code{Pat-1981-WB3/BX200} occupe
ainsi cinq lignes : un essai de traction et quatre essais Charpy, de \SI{39}{\joule} à
\SI{-60}{\celsius} à \SI{114}{\joule} à \SI{0}{\celsius}. On compte environ 620 cordons
pour 1652 lignes. Nous gardons une ligne par essai pour conserver les températures d'essai
et les traitements comme entrées. Les lignes d'un même cordon étant très proches, il
faudra en tenir compte dans l'évaluation (section~\ref{sec:methodes}).

\outrouver{\code{context/columns.json} ; notebook, « Analyses complémentaires › Une ligne
correspond à un essai ».}

\subsubsection{Valeurs manquantes}

\paragraph{Dans les entrées.} De 0\,\% (carbone, manganèse) à 96\,\% (tungstène) de
valeurs manquent selon les colonnes (figure~\ref{fig:manquants}) ; aucune ligne n'est
complète. Ces absences dépendent de l'article : 21 des 46 articles ne donnent jamais le
nickel, et les lignes sans nickel sont à 88\,\% des soudures manuelles, contre 49\,\%
pour les autres. Les valeurs ne manquent donc pas au hasard (pas \emph{MCAR}) ; leur
absence s'explique par l'article, connu (\emph{MAR}). Un nickel omis parce que très
faible relèverait du \emph{MNAR}, possible mais invérifiable.

\paragraph{Dans les résultats d'essais.} Chaque ligne ne contient qu'un essai : chaque
grandeur n'est connue que sur une partie des lignes (environ 880 pour le Charpy, 700 à 780
pour la traction). Ces cases sont structurellement vides, pas manquantes : l'essai n'a
pas été fait sur cette éprouvette. Fusionner les lignes d'un même métal ne les remplirait
pas sans perte, car un métal est souvent testé en Charpy à plusieurs températures : une
seule ligne ne garderait qu'un de ces essais, ou leur moyenne, sans sens physique. Ces
cases ne sont donc jamais complétées, et chaque modèle n'apprend que sur les articles qui
ont fait son essai (26 sur 46 pour le Charpy). La dureté (138 lignes), la FATT (31) et la microstructure
(environ 90) sont trop rares pour être apprises.

\outrouver{notebook, « Analyses complémentaires › Valeurs manquantes des entrées » et
« Lignes complètes ».}

\subsubsection{Types de variables}

\begin{itemize}
  \item \textbf{Numériques} : composition, réglages de soudage, traitement et résultats
    d'essais. Certaines valeurs sont écrites en texte (\code{<0.002}, \code{67tot33res},
    \code{150-200}) et doivent être converties (section~\ref{sec:preprocessing}).
  \item \textbf{Catégorielles nominales} : type de courant, polarité et procédé. Aucune
    n'est ordinale : la polarité \code{0} ne se situe pas entre \code{+} et \code{-},
    elle désigne le courant alternatif. Les catégories sont très déséquilibrées : soudage
    manuel MMA sur 1140 lignes sur 1652, courant alternatif sur 42 seulement.
\end{itemize}

\outrouver{notebook, « Data Exploration » (valeurs non numériques, fréquences des
catégories).}

\subsubsection{Unités de mesure}

La composition est donnée en \% pour les éléments principaux et en
ppm pour les traces. Ce double usage est une source d'erreurs : des valeurs absurdes dans
l'unité de leur colonne ont révélé des ppm notés en \%, et l'inverse
(tableau~\ref{tab:decodage}). Une fois corrigé, il est sans effet sur les modèles après
standardisation, mais une unité commune sera nécessaire pour les recommandations.

\afaire{\item Piste à vérifier en bibliographie : l'effet d'un traitement thermique
    dépend d'une combinaison de sa température (en kelvins) et du logarithme de sa durée
    (paramètre de Hollomon-Jaffe). Créer cette variable pourrait aider les modèles.}

\subsubsection{Distributions et valeurs extrêmes}

\paragraph{Grandeurs à prédire.} Les propriétés de traction ont des distributions
régulières (figure~\ref{fig:cibles}). Le Charpy présente deux pics : 356 lignes à
exactement \SI{100}{\joule} et 118 à \SI{28}{\joule}. Certains articles, surtout ceux
d'Evans, ne donnent que la température à laquelle la soudure atteint ces seuils :
« \SI{-38}{\celsius}, \SI{100}{\joule} » signifie « \SI{100}{\joule} atteints à
\SI{-38}{\celsius} ». Ce sont des points de la courbe énergie--température, calculés et
non mesurés.

\begin{figure}[!htbp]
  \centering
  \includegraphics[width=\linewidth]{cibles}
  \caption{Distribution des cinq grandeurs à prédire ($n$ : nombre de lignes où elle est
    mesurée).}
  \label{fig:cibles}
\end{figure}

\paragraph{Valeurs extrêmes.} La règle de l'IQR ne convient pas : 64\,\% des
températures entre passes valent exactement \SI{200}{\celsius}, et elle déclarerait
aberrantes toutes les autres. Nous avons donc examiné les valeurs extrêmes une à une.
Ce sont presque toujours des mesures réelles, que nous gardons :
\begin{itemize}
  \item des séries volontaires (10 soudures très riches en soufre et en phosphore) ;
  \item des soudures hors norme mais cohérentes (une soudure à 9\,\% de chrome a la plus
    forte résistance et le plus faible allongement) ;
  \item un effet de source (l'article le plus ancien, Wolst-1974, détient quatre valeurs
    extrêmes).
\end{itemize}

\outrouver{revue détaillée : \code{preprocessing/outlier_*.py}, conclusions dans
\code{context/columns.json} (champ \code{outliers}).}

\subsubsection{Relations entre variables}

\paragraph{Corrélations.} La figure~\ref{fig:correlations} donne la corrélation de
Spearman (sur les rangs, donc peu sensible aux valeurs extrêmes) entre entrées et
grandeurs à prédire. Elle indique un lien, pas une cause.
\begin{itemize}
  \item Les éléments d'alliage (Cr, Mo, V, Nb) vont avec plus de résistance et moins de
    ductilité.
  \item Les réglages de soudage (courant, tension, énergie) et les impuretés (S, P) vont
    avec moins de striction.
  \item Le Charpy n'est fortement lié à aucune entrée de composition ou de soudage ; sa
    variable la plus liée est la température d'essai ($r = 0{,}46$).
\end{itemize}

\begin{figure}[!htbp]
  \centering
  \includegraphics[width=0.62\linewidth]{correlations}
  \caption{Corrélation de Spearman entre entrées numériques et grandeurs à prédire. Les
    valeurs $|r| \geq 0{,}4$ sont écrites.}
  \label{fig:correlations}
\end{figure}

\paragraph{Analyse en composantes principales (ACP).} Appliquée aux 995 métaux distincts
(un métal testé plusieurs fois compte une fois), après imputation par la médiane et
standardisation (figure~\ref{fig:acp}) :
\begin{itemize}
  \item \textbf{les entrées ne sont pas redondantes} : deux axes n'expliquent que 38\,\%
    de la variance ; nous les gardons toutes ;
  \item \textbf{l'axe 1 (21\,\%) sépare aciers alliés et ordinaires} : les soudures
    riches en Cr, Mo, V, N et Nb sont presque toujours traitées thermiquement (94\,\% de
    celles à plus de 1\,\% de chrome), donc ces variables varient ensemble. On ne pourra
    pas distinguer l'effet des alliages de celui du traitement ;
  \item \textbf{l'axe 2 (17\,\%) correspond à la chaleur apportée par le soudage} : il
    sépare l'arc submergé (\SI{640}{\ampere} médian) de l'arc manuel (\SI{170}{\ampere}).
    Courant, tension et énergie de soudage varient ensemble, car l'énergie se calcule à
    partir des deux autres et de la vitesse : leurs effets ne pourront pas être séparés.
    Les recommandations porteront sur l'énergie, toujours renseignée et qui intègre la
    vitesse.
\end{itemize}
Ces résultats restent les mêmes avec d'autres méthodes d'imputation. Les modèles devront
couvrir des familles de soudures très différentes.

\begin{figure}[!htbp]
  \centering
  \includegraphics[width=\linewidth]{acp}
  \caption{ACP des entrées numériques : (a) variance expliquée par composante,
    (b) cercle des corrélations (flèches bleues : variables bien représentées), (c)
    projection des 995 métaux, par famille de procédé.}
  \label{fig:acp}
\end{figure}

\outrouver{notebook, « Analyse descriptive » (figures enregistrées dans
\code{report/figures/}) et « Énergie de soudage, courant et tension ».}

\afaire{\item Ajouter au notebook la comparaison des méthodes d'imputation avant l'ACP
    (médiane, plus proches voisins, imputation itérative).
  \item À décider en équipe : le traitement à \SI{250}{\celsius} / \SI{14}{\hour}
    (333 lignes, surtout de traction) est-il un simple dégazage équivalent à « brut » ?
  \item Recopier les résultats de traction sur les lignes Charpy du même état, utile
    seulement si la qualité combine résistance et ténacité.}
